\section{Transverse Fourier products and geometric transfer}
\label{sec:transfer}

We continue with the induced finite-module realization of
Section~\ref{sec:links}. All target indices in this section are indices
of a measured copy of the source lattice $\Gamma$. The comparison group
$\Lambda$ is retained in the coefficient algebra. In particular, none of
the flag spaces used below is assigned an action of $\Lambda$.

Our objective is a sharp flag measurement at each source flag, together
with an exact equation for the residual projectivity on a missing panel.
The Fourier product estimate first compares vertex paths and identifies
their target types. Isolated types then require a separation argument for
angular forms; in higher-rank components, incidence range containments
propagate sharpness through the building. The resulting ordered gallery
identities supply the residue equation used by the local panel arguments.

\subsection{A transverse product estimate}

We first state the estimate with its trace dependence visible. A finite
coefficient algebra has a fixed faithful finite trace. Matrix sizes are
fixed during every limit. We may normalize the coefficient trace for the
proof; returning to a trace of mass $t$ multiplies an operator-norm-only
$L^2$ bound by $\sqrt t$. More generally, the homogeneous estimate below
uses the actual Hilbert norms and does not require a normalization.
Rectangular matrices are included by embedding them in one fixed square
matrix algebra and using the corresponding corners.

For a scalar function $f$ of the measured Fourier index, let $D_f$
denote diagonal multiplication of the Fourier array. It is an $L^2$
operator, and we do not assume that it preserves operator-bounded
columns. Consider
\begin{equation}
 (D_g B_0)(D_f A_0),\qquad\text{with index product }qp.
 \label{eq:input-2}
\end{equation}
Here $A_0,B_0$ are operator bounded, $p$ is an index of $A_0$, and $q$
is an index of $B_0$. Write their Fourier expansions as
$A_0=\sum_p a_pw_p$ and $B_0=\sum_q b_qw_q$, with
operator-valued coefficients in the unmeasured algebra.

\begin{proposition}[Transverse products]
\label{prop:transfer-transverse}
Fix face types, a thickness bound, and sufficiently small compact
transverse face patches at the basepoint $o$. Suppose the supports of
$f$ and $g$ lie within that thickness of the respective positive cones,
with all specified positive gaps sufficiently large, and that the
ideal directions of $po$ and $q^{-1}o$ belong to the fixed transverse
patches. Then
\[
 \mu(qp)=\mu(q)+\mu(p)+O(1),
\]
where the error is uniform on the supports. The product in
\eqref{eq:input-2} has a uniform tracial $L^2$ bound in either of the
following situations:
\begin{enumerate}[label=\textup{(\roman*)}]
\item For each coordinate contributed by both cones, one of its two
supports is restricted to a moving interval of a fixed bounded width.
The separate cutoff factors may be arbitrary uniformly bounded
functions.
\item One of the face types is a single vertex. The cutoff factors are
smooth with uniformly bounded sufficiently high seminorms in fixed
compact normalized matrix charts; at finite places this means locally
constant with a uniformly bounded level.
\end{enumerate}
To specify the retained cutoffs in the estimate, write
$f=h_Af_1$ and $g=h_Bg_1$. In case~\textup{(i)}, these four
factors may be arbitrary uniformly bounded scalar functions. In
case~\textup{(ii)}, each factor has the smooth seminorm bounds or
fixed clopen level just specified. Then
\begin{equation}
 \big\|(D_gB_0)(D_fA_0)\big\|_2^2
 \leq C\|A_0\|_\infty\|B_0\|_\infty
       \|D_{h_A}A_0\|_2\|D_{h_B}B_0\|_2.
 \label{eq:transfer-homogeneous}
\end{equation}
The constant includes the stated bounds on the scalar factors.
The same estimate holds for a joint symbol, meaning the bilinear
Fourier sum
\[
 \sum_{q,p}h_B(q)h_A(p)m(q,p)b_qa_p\,w_qw_p,
\]
where $w_qw_p$ is the actual twisted monomial product and $m$ is a
smooth joint function on the fixed compact normalized charts, with
sufficiently many uniformly bounded derivatives. At finite places use
a fixed clopen level. The support and transverse restrictions remain
in force. Consequently an
additional smooth cutoff factor whose one-column Hilbert norm tends
to zero makes the product tend to zero, provided the other bounds and
the other cut Hilbert norm remain bounded. In particular, for a fixed
finite trace of mass $t$,
\[
 \big\|(D_gB_0)(D_fA_0)\big\|_2
 \leq C\sqrt t\,\|A_0\|_\infty\|B_0\|_\infty.
\]
\end{proposition}

\begin{proof}
We separate the geometric and Fourier parts of the argument.

\emph{Common apartments and addition.}
On compact transverse patches, the two ideal faces admit common
apartment or common flat sections whose origins remain in a bounded
set. After erasing the bounded thickness, write the two endpoints in
such coordinates as
\[
 po=z+A+O(1),\qquad q^{-1}o=z-B+O(1),
\]
with $A,B$ in the specified positive faces and $z$ uniformly bounded.
The vector distance between these endpoints is $A+B+O(1)$.
The vector-distance inequality under bounded endpoint perturbations
therefore proves the Cartan addition assertion. This argument is made
in the whole ambient product; it uses no discreteness of a proper
projection of the lattice.

We need the following stability of interpolation. From a fixed
endpoint, multiply the simple coordinates $L_i$ of a dominant vector
by arbitrary fractions $t_i\in[0,1]$. The resulting interpolation is
uniformly Lipschitz in both endpoints. In an archimedean polar chart,
the angular differential in a positive-root direction is multiplied by
\[
 \frac{\sinh\bigl(\alpha(t\cdot L)\bigr)}
      {\sinh\bigl(\alpha(L)\bigr)}\leq 1,
\]
and the radial differential is the uniformly bounded linear map
$L\mapsto t\cdot L$. A wall stabilizer fixes the interpolated vector
because zero simple coordinates remain zero. Thus the map extends
continuously across the walls; approximation of paths by regular paths
gives the same Lipschitz bound there. Reversing the endpoints uses
opposition and the complementary fractions.

For a finite-place building, the corresponding apartment-independent
interpolation and its uniform Lipschitz bounds in both endpoints are
Lemma~\ref{lem:setup-building-interpolation}. The product of these
real and finite-place bounds gives the assertion for the whole ambient
space.

\emph{The common shift.}
Suppose $qp=q'p'$ and set $k=p(p')^{-1}$. Then $k$ sends the primed
endpoint pair to the unprimed pair. Write the primed coordinates as
$z'+A',z'-B'$, with the same bounded errors. Cartan addition gives
\[
 A+B=A'+B'+O(1).
\]
Interpolate from the negative endpoint using the fractions
\[
 t_i=\frac{B_i'}{A_i'+B_i'},
\]
with any choice in $[0,1]$ when the denominator vanishes. These
fractions return $z'$ in the primed model. The Lipschitz bound and
equivariance of interpolation yield
\begin{equation}
 ko=z+B'-B+O(1).
 \label{eq:transfer-common-shift}
\end{equation}
There is no small-denominator loss: the fractions are bounded and are
held fixed when applying the Lipschitz estimate.

In case~\textup{(i)}, every overlapping coordinate has bounded variation
on at least one side. The relation between $A+B$ and $A'+B'$ then bounds
the corresponding coordinate of $B'-B$. Outside the overlap, one side
is bounded by the thickness, so those coordinates are bounded as well.
Thus $k$ belongs to a fixed compact subset of the whole product $G$.
Its intersection with $\Gamma$ is finite, uniformly in the locations
of the moving width intervals.

In case~\textup{(ii)}, an unbounded difference can occur only in
the single overlap coordinate $i$; if there is no overlap we are
already in the bounded-shift case. Put $\delta=B'_i-B_i$.
All other coordinates of $B'-B$ are bounded, and
$A_i-A'_i=\delta+O(1)$. The collision identity gives
\[
 p=kp',\qquad q'=qk,\qquad (q')^{-1}=k^{-1}q^{-1}.
\]
Apply Equation~\eqref{eq:transfer-common-shift} also to the reversed
pair. In the unprimed and primed apartments, respectively, it gives
\[
 ko=z+\delta e_i+O(1),\qquad
 k^{-1}o=z'-\delta e_i+O(1).
\]
The origins and both errors are uniformly bounded. Let $U$ be the
patch of type-$i$ vertices incident to the positive endpoint patch,
and $V$ the patch of type-$i^*$ vertices incident to the negative
endpoint patch. Choose buffered larger patches $U^+,V^+$ whose
every pair is opposite. For sufficiently large positive $\delta$,
the forward vertex of $k$ lies in $U^+$ and its reverse vertex,
the forward vertex of $k^{-1}$, lies in $V^+$. These conclusions
follow uniformly from the bounded-distance ray formulas. For large
negative $\delta$ the memberships reverse. The larger patches can
be chosen disjoint in the visual boundary, so a fixed sufficiently
large $k$ cannot have both orientations. Bounded $\delta$ gives
bounded $k$, already treated above. The negative ray has the opposed
type $i^*$; the two types need not coincide.

In the positive orientation the shorter arguments are $p'$ and
$q^{-1}$, and the length-increasing substitutions are
\[
 p'\longmapsto kp',\qquad
 q^{-1}\longmapsto k^{-1}q^{-1}.
\]
The reverse vertex of the first multiplier is in $V^+$, transverse
to the entire buffered positive input patch. The reverse vertex of
the second is in $U^+$, transverse to the entire negative input
patch. Thus the witness for $k$ gives uniform transversality on
both whole input patches. The negative orientation interchanges the
primed and unprimed indices and replaces $k$ by $k^{-1}$.

For clarity, let $R$ be one of the fixed finite-dimensional
representations and put $S(g)=R(g)/\|R(g)\|$. On either substitution $r\mapsto ur$, the
transverse Cartan-addition estimate applies to every $r$ in the
buffered initial patch. Highest-weight norm comparison gives
\[
 \|S(u)S(r)\|\geq c>0,\qquad
 \|S(r^{-1})S(u^{-1})\|\geq c>0,
\]
uniformly in the supported $u,r$. The second bound applies
opposition to Cartan addition and uses $(ur)^{-1}=r^{-1}u^{-1}$.
The inverse coordinates are the separately normalized matrices of
$g^{-1}$, so no inverse of a possibly singular limiting matrix is
being taken. There are finitely many chosen modules, and the buffer
preserves a common lower bound on fixed larger coordinate boxes.
Normalized multiplication consequently has uniformly bounded
smooth derivatives on those boxes. At a finite place, maximal-coordinate
ratio charts have uniformly nonzero denominators; the analytic
substitutions pull a fixed clopen test back to a common bounded level.

\emph{Gram coefficients and twists.}
Expand the squared norm in \eqref{eq:input-2} and group the pairs of
indices with $qp=q'p'$ by $k=p(p')^{-1}$. After trace cyclicity, the
summand is the pairing of two Gram coefficients, one made from
$a_{p'}a_p^*$ and the other from $b_q^*b_{q'}$, with adjoints changed
if the opposite orientation is used. Here the lower-case letters are
the operator-valued Fourier coefficients. The cocycle factors are
exactly the factors in these two actual adjoint products: this follows
by associativity of the twisted Fourier monomials, or directly by the
scalar cocycle identity. In particular, we do not replace the
operator-valued convolution by an absolute-value convolution.

For the finitely many bounded shifts, row and column
Cauchy--Schwarz bound each weighted Gram coefficient in operator norm
by the appropriate uncut column norm squared. Its trace norm is
bounded by the relevant cut Hilbert norm squared. Combining the two
bounds symmetrically and then applying tracial Cauchy--Schwarz gives
\eqref{eq:transfer-homogeneous}, with a factor equal to the uniformly
bounded number of shifts. Arbitrary separate bounded cutoffs cause
no problem in this part.

For the unbounded shifts, use the orientation just described.
When $p=kp'$ is the length-increased side, the first Gram weight is
\[
 f(kp')\overline{f(p')}
 =\overline{h_A(p')}\,
       f(kp')\overline{f_1(p')}.
\]
Retain $\overline{h_A(p')}$ and expand the remaining factor in
the shorter argument $p'$. On the other side, retain $h_B(q)$ in
\[
 g(q)\overline{g(qk)}
 =h_B(q)\,g_1(q)\overline{g(qk)};
\]
its shorter geometric argument is $q^{-1}$. Conjugating a Gram
coefficient may conjugate these formulas, with no effect on the cut
Hilbert norm. In the opposite orientation interchange primes, retaining
$h_A(p)$ and $h_B(q')$. The longer-argument occurrence of each retained
factor remains inside the expanded expression; its uniform substitution
bounds are therefore still required in case~\textup{(ii)}.

Extend each remaining expression using an initial-patch cutoff with
slack. For the first Gram sequence, choose conjugated basis functions
so that its full weight has the form
\begin{equation}
 w_k(r)=\sum_n c_n(k)\overline{h(r)u_n(r)},\qquad
 \|u_n\|_\infty\leq1,\qquad
 \sum_n\sup_k|c_n(k)|<\infty.
 \label{eq:transfer-unary-expansion}
\end{equation}
The last, stronger, order of sum and supremum is essential. It follows
from Fourier expansion on a slightly larger real coordinate box and
uniform bounds on sufficiently many derivatives. At finite places
there is instead a finite basis of fixed-level clopen masks. Products
of these constructions cover mixed charts. The other Gram sequence
has the same form after the indicated adjunction.

For example, when $p=kp'$ is the length-increased side, $h=h_A$
and the first Gram
sequence is a sum of coefficients of
$A_0(D_{hu_n}A_0)^*$ with scalar weights $c_n(k)$. Minkowski,
Plancherel, and the Hilbert ideal inequality give
\[
 \left\|\left(\sum_n c_n(k)
      [A_0(D_{hu_n}A_0)^*]_k\right)_k\right\|_{\ell^2(L^2)}
 \leq
 \left(\sum_n\sup_k|c_n(k)|\right)
 \|A_0\|_\infty\|D_hA_0\|_2.
\]
On the other side, use the shorter of $q^{-1},(q')^{-1}$ to obtain
the analogous bound for $B_0^*B_0$. If $B'-B$ is large and positive,
the initial arguments are $p'$ and $q^{-1}$; the negative orientation
is symmetric. Cauchy--Schwarz in the shift variable proves
\eqref{eq:transfer-homogeneous}. The estimate also proves the
small-factor assertion; a square-root rate for the product norm is
sufficient.

Smooth joint symbols are handled by summably separated expansions,
retaining the thickness and transverse restrictions with slack in
every term. Finally approximate general bounded columns by uniformly
bounded Fourier polynomials. Unary products converge in $L^1$, and
lower semicontinuity gives the asserted $L^2$ bound. No step requires
operator-norm control of an angularly cut column.
\end{proof}

\begin{remark}[Uniformity used below]
\label{rem:transfer-uniformity}
In case~\textup{(i)}, the constant is insensitive to the locations of
the bounded-width intervals and to arbitrary bounded separate masks.
In case~\textup{(ii)}, all smooth seminorms and all finite-place clopen
levels must be held fixed during a use of the estimate. A gap threshold
may be fixed first and exhausted only after the relevant inner limit.
For the fixed matrix amplifications below, the cut Hilbert norms in
\eqref{eq:transfer-homogeneous} are uniformly bounded. Thus no estimate
independent of an unbounded unnormalized matrix trace is needed.
\end{remark}

\subsection{Fullness and the permutation of types}

For an archimedean target vertex module $V$, let $n(v)$ be a unit
representative of its highest line; the rank-one projector
$n(v)n(v)^*$ is independent of the choice of phase. We also allow the
trivial one-dimensional module. At a finite target place only this
trivial-module version is used in the next lemma.

\begin{lemma}[Fullness of open patches]
\label{lem:transfer-fullness}
Let $\Phi_z$ be a nonzero family of image measurements constructed in
Proposition~\ref{prop:links-link}, transported over its source label
space and normalized by its positive scalar stratum mass to be unital. For every
nonempty open target patch $O$, the closed sum, over conull source
labels $z$, of the ranges of
\begin{equation}
 \int_O n(v)n(v)^*\otimes\Phi_z(dv)
 \label{eq:transfer-line-effect}
\end{equation}
is the full amplified input space. For the trivial module this is the
ordinary patch-fullness assertion. The assertion passes to the normal
copies of the measurements in the three-slot system.
\end{lemma}

\begin{proof}
First omit $O$. Source covariance makes the uncut range join a
projection over the finite algebra $(\mathcal A_1)^\sigma$. Target
covariance transports its range by the finite-dimensional linear
target matrices; positive changes in the normalization of $n(v)$ do
not change that range. Induce the target lattice action, with its
finite invariant trace, to the ambient group. The induced range field
then intertwines a trace-isometric ambient action with the inverse
linear transport on $V$.

We spell out the finite-module fact used here. Let a positive matrix
$A$ have weight spaces $V_1,\ldots,V_s$ with strictly increasing
eigenvalues, and let $P$ be a projection over a finite algebra. Filter
$\ran P$ by its intersections with $V_1\oplus\cdots\oplus V_j$.
Take the closures of the leading projections of these intersections
onto $V_j$. Polar decomposition shows that their dimensions are the
successive differences of the filtration dimensions, so their
orthogonal sum $Q$ has the same trace as $P$. If a vector belongs to
the $j$th filtered intersection, its image under $A^n$, divided by the
$j$th eigenvalue to the power $n$, converges to its leading projection.
Consequently vectors in $\ran Q$ are strongly approximated by vectors
in $A^n\ran P$. Since the latter range projections all have the same
finite trace, they converge to $Q$ in trace $L^2$, equivalently in
trace norm for projections.

Apply this pointwise to the induced field and integrate, using the
finite trace. A convergent orbit of a trace-isometric action is fixed:
the distance between two consecutive orbit points is constant and
also tends to zero. The field is therefore invariant under every
positive split matrix and its conjugates. These generate the relevant
connected simple group. Irreducibility of the complexified real
module forces the range to split off the whole matrix factor. The
remaining projection descends to the scalar joint fixed algebra, so
the uncut join is full. For an amplified real module in a product,
use the matrix representation only in that archimedean component and
the trivial representation in all others; the same argument applies.

Now restore $O$. The finite trace dimension of the patch join is
constant along target lattice translates, because the transporting
linear matrices are invertible. Given any finite scalar test measure
on the target flags, choose a regular contracting direction whose
attractor is in $O$ and whose repeller is transverse to almost every
point of that measure. Such a repeller exists by the nullity of the
proper incidence loci. All-factor mixing on the lattice quotient
supplies lattice elements whose two compact frames approximate the
chosen frames along increasingly long regular translations: two
fixed open frame neighborhoods have positive translated intersection
in the quotient, and the neighborhoods may subsequently be shrunk.
Contraction away from the repeller implies that suitable translates
of $O$ exhaust the test measure.

Apply this to a faithful scalarization of a countable total family of
label-averaged effects. The complements of the translated patch joins
then have trace tending to zero. One way to see the last implication
is to take a faithful positive element generated by that total family
and truncate it spectrally away from zero. The constant trace
dimension of the translated joins must therefore be the full trace.
This proves patch fullness. The normal-copy assertion follows from
the normal ampliation statement in Proposition~\ref{prop:links-link}.
\end{proof}

A normalized image branch of a pure vertex path will be called a
\emph{candidate}. Its label $z$ is a source vertex, whereas its values
are target vertices. Different candidate choices initially include
different paths, limits, and positive strata. Each such choice at a
fixed source vertex gives a transported family by
Proposition~\ref{prop:links-link}. Face-parabolic invariance makes
this transport well defined on the vertex space. Strong continuity of
$U$ and local sections of the homogeneous quotient then make its
effects strongly continuous in the source label for each fixed target
test. This continuity is used only for these constructed families.

\begin{proposition}[Type comparison]
\label{prop:transfer-types}
There is a unique permutation $p$ of the full relative type list
$\Pi$ such that every candidate at source type $i$ takes values in
target type $p(i)$. It preserves opposition:
$p(i^*)=p(i)^*$. For noncoaxial source vertices of dual types, all
ordered products of candidates vanish on transverse target patches.
The permutation preserves the collection of isolated types.
\end{proposition}

\begin{proof}
Choose paths $t,a$ whose inverses tend to the two source vertices in a
common apartment, using parabolic invariance to choose the split
frames. Take successive limits, with $t$ faster. Pairing their ordered
cut effects between bounded columns and using trace cyclicity gives
a pairing of
\[
 (D_AX_t f)(D_BX_a f')^*
\]
with the quotient-path matrix $X_tX_a^*$, after the appropriate base
reindexing. The latter is the transporter of $ta^{-1}$ up to a common
scalar base phase. Its index product is $pq^{-1}$, and the intermediate
ideal directions are $p^{-1}o,q^{-1}o$, in accordance with the inverse
Fourier convention.

If the source vertices are coaxial, the quotient path has exactly
one target gap by Proposition~\ref{prop:links-gap-count}. Nondual
transverse target types would force two distinct large gaps in the
cut product. The bounded-overlap version of
Proposition~\ref{prop:transfer-transverse} bounds its norm uniformly
at fixed angular patches and thickness, independently of the large-gap
thresholds. The localized norm of the uncut quotient matrix tends
to zero on the forbidden gap strata; this is checked on the finite
frame columns. Cauchy--Schwarz therefore makes the pairing zero.

If the source vertices are noncoaxial, the quotient has at least two
gaps. Dual transverse target vertex cones put the cut product in a
fixed thickened single-vertex cone. The smooth version of
Proposition~\ref{prop:transfer-transverse}, with fixed smooth cuts,
again compares it with a forbidden localized quotient norm and gives
zero. Fix the smooth thickness and angular cuts first, then exhaust
the large-gap and thickness windows. This proves the stated ordered
vanishing without a claim about uncut operator norms.

Fullness supplies a nonzero ordered cut product for generic source
pairs: successively apply Lemma~\ref{lem:transfer-fullness} to the two
families. Generic dual source vertices are opposite and hence
coaxial, whereas generic nondual source types are noncoaxial. Thus
any candidate choices at dual source types have dual target types,
and choices at nondual source types have nondual target types.
If $T_i$ is the set of possible target types at source type $i$, every
$j\in T_i$ and $k\in T_{i^*}$ satisfy $k=j^*$. Since the sets are
nonempty, each $T_i$ is a singleton. Preservation of both duality and
nonduality makes the resulting map injective, and the finite common
type list makes it a permutation.

At equality of higher-rank self-opposite source vertices, the
nonopposition exclusion extends by continuity. Indeed such a pair is
a limit of distinct nonopposite pairs, obtained by varying a chamber
across an adjacent panel in the connected diagram. If the image type
were isolated, nonopposition would mean equality. Diagonal support of
two unital effects would then make them sharp and equal. The same
argument applies along the adjacent chamber changes, and gallery
connectivity would make the candidate family constant. It would lie
in the finite source-fixed algebra and be target equivariant. Taking
its invariant trace would give an invariant target flag probability,
contradicting flag contraction as in Lemma~\ref{lem:transfer-fullness}.
Thus a higher-rank self-opposite type cannot map to an isolated type.
A non-self-opposite type cannot do so either, by opposition
preservation. Since $p$ is a permutation, it permutes the isolated
types.
\end{proof}

\subsection{Equality at isolated types}

All isolated types under consideration are archimedean. Fix one such
source type and its matched target boundary $Z$. The sharpness argument will produce a scalar functional bounded by
the angular norms at two independently moved source labels. We first
show that such a nonzero functional cannot exist. The proof encodes
equivariant positive maps by frame functionals, uses double ergodicity
to separate them, and returns from dominated maps to angular forms.
Neither candidate family is assumed sharp.

\begin{lemma}[Disjoint angular forms]
\label{lem:transfer-disjointness}
For almost every independent pair of source labels, possibly taken
from two different candidate families at the fixed isolated type,
there is no common nonzero target-equivariant positive map dominated
by the two candidate maps, where domination is complete-positive
domination. After induction, there is no common nonzero linear or
conjugate-linear functional on
$C^\infty(Z)\odot\mathcal H$ bounded by both angular form norms
\begin{equation}
 b\longmapsto
 \left(\int_Z\langle b(v),\Phi_x(dv)b(v)\rangle\right)^{1/2}.
 \label{eq:input-3}
\end{equation}
The generic conclusion holds throughout the transverse source orbit
by source covariance.
\end{lemma}

\begin{proof}
Undo source induction. Simultaneous covariance of the base and source
label yields measurable unital completely positive maps
\[
 F_x:C(Z)\longrightarrow B(H),\qquad
 F_{\gamma x}=\pi_0(\gamma)F_x\pi_0(\gamma)^*.
\]
They are equivariant for $R$, where the designated measured Gamma
factor acts on $Z$ and the other factors act trivially. The assertion
follows either from the lattice cocycle formula or from Haar
disintegration on the source homogeneous space. Scalar projective
phases disappear in conjugation. The source lattice action on pairs
of this boundary is ergodic: the open-pair stabilizer contains a split
subgroup escaping in the active factor and all complementary ambient
factors, so all-factor mixing on the finite-volume quotient applies.

Let $\mathfrak C$ be the separable matrix-amplified full twisted
crossed product associated with $C(Z)$ and $R$. Encode $F_x$ by its
minimal covariant Stinespring dilation~\cite{Stinespring1955} and by the vector obtained from the
embedded columns of $e$, arranged as the matrix $e$. The resulting
positive functional $\phi_x$ on $\mathfrak C$ has norm $d$.
The group and matrix operators act on these embedded columns by right
multiplication. The regular-normal central summand of the group
algebra bidual, mapped normally into the crossed-product bidual,
therefore supplies canonical right multipliers by matrices over the
regular group von Neumann algebra. The vector functional is supported
on the corresponding right-$e$ corner. This support assertion concerns
the embedded matrix frame; angular multiples generating the dilation
need not themselves have regular group support.

Source translation acts on $\phi_x$ by an inner automorphism of this
fixed corner. Indeed, in the vector formula for the translated map,
$e$ is replaced by
\[
 \pi_0(\gamma)^*e=e\pi_0(\gamma)^*.
\]
The right side uses the same canonical right corner unitary for every
label and every dilation. Equality on crossed-product polynomials
extends normally to the bidual. Extending that corner unitary by the
identity on its complementary projection, if desired, gives a unitary
of the bidual. In particular every central summand is preserved by
these inner automorphisms.

The positive subfunctionals of $\phi_x$ encode exactly the
corresponding equivariant subordinate maps. To verify this, note that
the matrix vector is cyclic: its group and matrix translates span the
embedded input, and applying angular functions generates the minimal
dilation. A dominated positive functional is represented by a
positive contraction in this cyclic commutant. It commutes with the
matrix index and group representation; compressing weighted angular
multiplication gives a subordinate equivariant map on the original
input. Conversely, a subordinate map gives the dominated form in the
same dilation. This is the covariant form of the CP domination
correspondence \cite[Theorem~1.4.2]{Arveson1969}. Its coefficients between group and matrix translates
of the cyclic vector give uniqueness. This correspondence also
intertwines source conjugation by the corner unitaries just described.

Fix the central support $z$ of one functional representation from the
second family. The summand $z\mathfrak C^{**}$ has separable predual,
because that cyclic representation acts on a separable Hilbert space.
The components $\phi_x^z$ form a measurable norm-separable family.
Here is the measurable-selection detail. In the Polish norm topology
of the positive normal ball of $z\mathfrak C^{**}$ of radius $d$,
the condition $0\leq\theta\leq\phi_x$ is Borel: test it on a
countable norm-dense family of positive elements of $\mathfrak C$.
Put $n(x)=\sup\{\|\theta\|:0\leq\theta\leq\phi_x,
\ \theta\text{ is }z\text{-normal}\}$. This function is upper
semianalytic. For each rational $q$, the Jankov--von Neumann selection
theorem gives, on $\{n>q\}$, a universally measurable dominated
normal functional $\theta_q(x)$ with norm greater than $q$
\cite{Kechris1995}. Choose measurable rational
$q_m(x)<n(x)$ with $n(x)-q_m(x)<1/m$, treating $n=0$ separately.
The maximal normal component $\phi_x^z$ dominates each such selection,
so positivity gives
\[
 \|\phi_x^z-\theta_{q_m(x)}(x)\|
 =n(x)-\|\theta_{q_m(x)}(x)\|<1/m.
\]
Thus $\phi_x^z$ is norm measurable on the completed boundary
probability space. No simultaneous measurable choice of the second
family's central supports is required.

The distance between $\phi_x^z$ and $\phi_y^z$ is invariant under the
diagonal source action, since both are conjugated by the same corner
unitary. Double ergodicity makes this distance essentially constant.
A measurable map into a separable metric space cannot have a strictly
positive constant distance at almost every pair: cover its range by
countably many balls of diameter smaller than that constant and take
one with positive preimage measure. Thus $\phi_x^z$ is constant almost
everywhere. If it were nonzero, the encoding would give a nonzero
source-invariant, target-equivariant positive map into the uninduced
finite complementary commutant. Its invariant finite trace would
produce an invariant flag probability, which is impossible by target
contraction. Hence $\phi_x^z=0$ almost everywhere.

For each fixed label of the second family, the two functionals are
therefore disjoint at almost every first label. Orthogonality is the
Borel condition that the distance equals the sum of the norms; the
norm is tested on a countable dense set in the separable algebra.
Fubini gives the first assertion for almost every independent pair.

Now reinduce and consider the two seminorms in
\eqref{eq:input-3}. If a nonzero scalar functional were bounded by both,
the parallel sum of the two squared forms would be nonzero. Construct
it by taking the direct sum of the two form completions and quotienting
by the closed span of pairs $(h,-h)$ coming from the common algebraic
domain. This subspace is reducing for smooth angular multiplication,
base multipliers, and the common target covariances, including their
adjoints. Angular multiplication descends, and compression at constant
input vectors gives a decomposable equivariant CP map dominated by
both original maps. It is nonzero: if every constant input vanished,
its angular multiples would vanish, although those vectors generate
the quotient. Disintegration gives forbidden subordinate maps on a
positive-measure set of base points. This contradiction proves the
form assertion. Finally, source covariance transports the generic
conclusion over the single transverse orbit.
\end{proof}

\begin{remark}[A consequence for sharp measurements]
\label{rem:transfer-disjoint-ranges}
If the two measurements in Lemma~\ref{lem:transfer-disjointness} are
sharp, the ranges $P_x(U)P_y(V)\mathcal H$, for disjoint angular
patches $U,V$, have dense sum. A vector orthogonal to all these ranges
would, by regular approximation of the spectral measures, have the
same spectral vector for the two PVMs. Pairing with that vector gives
a common functional bounded by the two forms, contrary to the lemma.
\end{remark}

\begin{proposition}[Sharpness at isolated types]
\label{prop:transfer-isolated-sharpness}
At an isolated source vertex, every candidate is sharp, and all
candidate choices agree.
\end{proposition}

\begin{proof}
Take fast and slow paths towards the same source vertex $x$. Suppose
that two angular effects on disjoint matched target patches have a
nonzero ordered product. Choose smooth buffered angular cutoffs with
bounded complementary thickness. Write $A$ for the fast cutoff and
$L$ for its one-ended limiting effect. Let $B_1$ be the slow cutoff,
used on slow-link columns, with its start patch disjoint from that of
$A$. The finite corner-frame identity and exhaustion of fixed windows
then give bounded columns $f,f',b'$ for which the following
conjugate-linear functional is nonzero. For
$b(v)=\sum_i h_i(v)b_i$, with smooth $h_i$ and bounded columns $b_i$,
it is the slow limit of
\begin{equation}
 \sum_i\left\langle
 D_{h_i(q^+)B_1}
       [b'(ay)^*X_a(y)f'(y)],\,
 b_i(ay)^*X_a(y)(Lf)(y)
 \right\rangle.
 \label{eq:input-4}
\end{equation}
Here the inner product is linear in its second argument and $q^+$ is
the image of the ordinary index $q$, that is, the reverse angular
variable in the matched component. Smooth boundary tests on the
large-gap windows are extended through the simple top singular line
in the normalized proximal matrix chart.

We show that this functional is bounded by the fast uncut form at
$x$. At fixed slow time, replace $Lf$ by its defining fast weak limit
along $t$. Trace cyclicity pairs
$X_tX_a^*b_i(ay)$ with the transverse product of the cut fast column
$D_AX_tf$ and the adjoint of the first column of
\eqref{eq:input-4}. Its indices are $p,q^{-1}$. Transfer $h_i$ to the
pairing column by evaluating it at the image of the inverse product
$qp^{-1}$, which approaches $q^+$.

The transfer error has a uniform quantitative form. Normalized
product formulas are smooth with bounded derivatives on the buffered
patches, by Proposition~\ref{prop:transfer-transverse}. Replace the
normalized matrix of $q$ by its rank-one top truncation. The two
angular tests agree exactly on that truncation. Smooth division, or
the first-order Taylor formula in a matrix chart, writes their
difference as a finite sum of entries of the rank-one defect times
symbols with uniformly bounded seminorms, with the original cutoffs
retained. The slow cut column multiplied by this defect tends to zero
in Hilbert norm, since the slow path has a unique vertex type and the
complementary thickness is fixed. The small-factor assertion of
Proposition~\ref{prop:transfer-transverse} makes the product error tend
to zero.

After transfer, the remaining cut product has a uniform Hilbert norm.
The squared norm of the summed pairing column has fast limit equal
to the fast uncut angular form on
$\sum_i h_i U_a^*b_i$. The split element $a$ fixes $x$, so parabolic
invariance identifies this with the form on $b$ at $x$. Extensions of
the transferred symbols away from the tested windows can be chosen
to have exactly the prescribed boundary values on every rank-one
limit. Cauchy--Schwarz therefore gives the desired form bound.

Independent contracted invariance at the slow reverse end permits
simultaneously replacing $b'$ and every $b_i$ by $U_n b'$ and
$U_n b_i$, for a reverse-radical element $n$. The functional itself
is unchanged. The bound just obtained is then the fast form on
$\sum_i h_iU_n b_i$, hence the form on $b$ at $n^{-1}x$; the slow
translation may be removed because it fixes $x$. These opposite
radical moves cover an open Bruhat cell. All constants may depend on
the fixed move, which is harmless. The form bounds extend to
arbitrary Hilbert coefficient columns by continuity. If the slow
path is nonmixed, $Lf$ is merely Hilbert, but every end column used
before the link limit is operator bounded; the one-bounded-input
extension in Lemma~\ref{lem:links-relative-boundedness} is exactly what is used.

Choose two generic moved labels in that cell. The same nonzero
functional is bounded by both of their forms, contradicting
Lemma~\ref{lem:transfer-disjointness}. Thus candidate products at the
same source vertex are supported on equality of target values.
For two unital POVMs $E,F$ with this property, regular approximation
gives $E(S)F(T)=0$ for disjoint Borel sets. Unitality then gives
$E(S)=E(S)F(S)=F(S)$ and $E(S)^2=E(S)$. Hence the candidates are equal
PVMs, as asserted.
\end{proof}

\subsection{Incidence range containments}

Let $F$ be a source face, let $z\sim F$ mean that $z$ is incident to
$F$, and fix a positive stratum of a mixed source face path towards $F$, or
of a pure vertex path $F$. Denote its normalized image-face marginal
by $\Psi_F$. The scalar mass of the unnormalized stratum is denoted
$c_F>0$. For an open patch $O$ of target type $p(i)$, let
\[
 \mathcal I_i(O)=\{x:\text{the target face }x
       \text{ admits an incident type-}p(i)\text{ vertex in }O\}.
\]
For an archimedean highest-line module $V$, put
\[
 R_x=\text{the orthogonal projector onto }
       \operatorname{span}\{n(v):v\sim x,\ \operatorname{type}(v)=p(i)\},
\]
and define
\begin{equation}
 L_z=\overline{\ran}\int n(v)n(v)^*\otimes\Phi_z(dv).
 \label{eq:transfer-Lz}
\end{equation}

\begin{proposition}[Incidence containment]
\label{prop:transfer-incidence}
The following inclusions hold:
\begin{enumerate}[label=\textup{(\roman*)}]
\item The support space of $\Psi_F(\mathcal I_i(O))$ is contained in
$\overline{\sum_{z\sim F}\ran\Phi_z(O)}$.
\item If the path is mixed and $p(i)$ is archimedean, in every even
highest power the support of
$\int R_x\otimes\Psi_F(dx)$ is contained in
$\overline{\sum_{z\sim F}L_z}$.
\end{enumerate}
In the proof of \textup{(i)}, a stronger individual-pair statement
holds. On a compact full-data window that sends a buffered transverse
end patch $V_1$ into $O$, if $f\in\ker\Phi_z(O)$ and
$b\in\overline{\ran}\Phi_z(V_1)$, then
$E(\mathrm{window})J(\bar b,f)=0$. On a mixed path $f$ may be any Hilbert column; on a nonmixed path
$f$ is required to be relatively bounded in the sense of
Lemma~\ref{lem:links-relative-boundedness}. In both cases $b$ may
be any vector in the stated closed end range.
\end{proposition}

\begin{proof}
We prove \textup{(i)} by annihilating the common kernel
\[
 \mathcal C=\bigcap_{z\sim F}\ker\Phi_z(O).
\]
First choose the domain on which this calculation will be made.  If the
path is mixed, take any $f\in\mathcal C$ and operator-bounded columns
$f_l\to f$ in Hilbert norm.  If it is a nonmixed vertex path supported
in $G_0$, put $N_0=\mathcal N^{G_0}$.  When $i$ lies in another
component, every type-$i$ label is incident to $F$, and ordinary
fullness gives $\mathcal C=0$, proving the inclusion.  Otherwise
Lemma~\ref{lem:links-relative-boundedness}\textup{(ii)} makes
$\mathcal C$ invariant under $N_0$.  Take $f$ in its dense relatively
bounded part and, by part~\textup{(iii)} of that lemma, choose
operator-bounded $f_l\to f$ satisfying
\[
 E_{N_0}(f_lf_l^*)\leq E_{N_0}(ff^*)\in N_0.
\]
The approximants need not belong to $\mathcal C$.  In either case,
for each $z\sim F$,
$\langle f_l,\Phi_z(O)f_l\rangle\to0$.

Fix a compact full-data window that sends a buffered transverse end
patch $V_1$ into $O$.  Such windows cover the image event in
\textup{(i)}.  Fix also $z\sim F$ and an operator-bounded column
$b_0$.  A vertex path $s^{-1}$ toward $z$, with a smooth angular cut
in $V_1$ and fixed complementary thickness, realizes its cut end
effect by an auxiliary weak limit.  Exhausting these cuts and
varying $b_0$ gives a total set in
$\overline{\ran}\Phi_z(V_1)$.

Keep the window, angular cuts, thickness, and bounded testing columns
$b'',f''$ fixed.  For each fixed main time $a$ and approximation
index $l$, insert the auxiliary $s$-limit.  Trace cyclicity pairs
\begin{equation}
 X_s(ay)X_a(y)f_l(y)
 \label{eq:transfer-incidence-pairing}
\end{equation}
with the product of cut versions of
\[
 X_s(ay)b_0(ay),\qquad
 b''(ay)^*X_a(y)f''(y).
\]
If the two product indices are $q,h$, their observed inverse product
is $h^{-1}q^{-1}$.  Transversality and the buffer put its pertinent
image in $O$ once the auxiliary gap is sufficiently large.
The effective norm of the pairing column in this inner limit is
therefore at most
\[
 \langle f_l,\Phi_z(O)f_l\rangle^{1/2}.
\]
The main split translation fixes $z$, because it acts trivially on
the Levi chamber refinements incident to $F$.  The bound is thus
independent of the main time.

Proposition~\ref{prop:transfer-transverse} gives a uniform Hilbert
bound for the other, cut product.  Its constant depends on the fixed
window, thickness, buffered patches, and the bounded columns
$b_0,b'',f''$.  It does not depend on $\|f_l\|_\infty$: $f_l$
appears only in the pairing column.  For each fixed auxiliary cut
and $b_0$, its limiting end vector is fixed.  Take the required
main-link limits with $l$ fixed, then let $l\to\infty$ against
that vector.  On a mixed path the tensor isometry justifies this
passage.  On a nonmixed path, the common conditional-expectation
bound on $f_l$ justifies it by
Lemma~\ref{lem:links-relative-boundedness}.  The displayed pairing
bound tends to zero.

Now keep $f$ fixed, exhaust the auxiliary thickness and angular
cuts, and take the closed span of the resulting end ranges.  The
same lemma gives continuity in the entire Hilbert end variable on
a nonmixed path; the mixed tensor isometry gives it on a mixed
path.  Consequently
\[
 E(\mathrm{window})J(\bar b,f)=0
       \qquad
       (b\in\overline{\ran}\Phi_z(V_1)).
\]
At the fixed-time auxiliary insertion the approximating input $f_l$
was operator bounded.  The subsequent Hilbert-input passages are
extensions of the limiting link; no finite-time product of two
arbitrary Hilbert columns is used.

For the individual-pair assertion, fix a single $z$ and repeat the
calculation with $f\in\ker\Phi_z(O)$.  In the nonmixed case take
$f$ relatively bounded.  Its bounded truncations still have the
same expectation bound and need not remain in the kernel; thus
this argument applies even when $i$ is outside the active component.
Kernel invariance was needed only for density in the range-inclusion
argument.  This proves the individual-pair assertion for every
Hilbert $f$ on a mixed path, and every relatively bounded $f$ on a
nonmixed path, with $b$ any vector in the stated closed end range.

Move the end occurrence also by the independent reverse radical.
As $z\sim F$ and these moves vary, the resulting labels cover a
transverse open cell, which is conull. Fullness and strong continuity
on fixed input vectors make these end ranges total. The link is
therefore annihilated on the window for every end vector.

We make the final frame step explicit, because the nonmixed case
cannot use tensor independence. First exhaust the compact full-data
windows to the event $A$ consisting only of the chosen gap stratum
and the image condition $\mathcal I_i(O)$. Monotone spectral
continuity gives $E(A)J(\bar b,f)=0$ for every end vector $b$.
This exhausted event is invariant under the ignored fixed left
Fourier shifts; the residual and reverse-window restrictions have
been removed. For the finite constant column frame $b_r$ of $e$,
the frame identity of Proposition~\ref{prop:links-link} is
\[
 c_F\langle f,\Psi_F(\mathcal I_i(O))f\rangle
   =\frac1d\sum_r\|E(A)J(\overline{b_r},f)\|^2=0.
\]
Divide by $c_F>0$, use positivity, and extend from the relatively
bounded common-kernel part by density. Orthogonal complements prove
\textup{(i)}. An arbitrary Borel output subprojection may be imposed
after the buffered ambient tests; no uniform estimate for a
discontinuous unexhausted window has been asserted.

For \textup{(ii)}, write $S(h)$ for the Hilbert-normalized real module
matrix at $h^{-1}$. If $(f_i)\perp L_z$ and $(b_j)$ belongs to the
amplified line range on a patch transverse to the output reverse
window, the same computation gives
\begin{equation}
 E(\mathrm{window})\sum_{ij}E(S_{ij}(h))J(\overline{b_j},f_i)=0.
 \label{eq:input-5}
\end{equation}
Indeed, the auxiliary fast symbol is $S(q)S(q)^*$ with its cutoffs,
and
\[
 S(h)S(q)S(q)^*=\rho(h,q)S(qh)S(q)^*,
\]
where the normalization ratio $\rho$ and its needed derivatives are
bounded on the transverse supports. Move $S(qh)^*$ to the full
pairing column. Its squared symbol tends in the auxiliary limit to
the line-projector effect defining $L_z$, with the same exact
stabilizer translation as above. It therefore vanishes on the
approximating kernel input in the limit.

The reverse radical may again move the end occurrence. Amplified
fullness now makes the resulting end columns span the whole amplified
end space. On a compact stratum window the surviving top block of
$S(h)$ has range containing every incident highest line, by
Lemma~\ref{lem:links-matrix-data}. Bounded local right inverses on its
constant-rank image, or finite local spanning columns followed by
output matrix multiplication, transfer the annihilation in
\eqref{eq:input-5} to the projector $R_x$. Compressing and using the
mixed input independence from Proposition~\ref{prop:links-link}
proves \textup{(ii)}. This amplified argument is used only for real
Hilbert modules; its finite-place replacement is given below.
\end{proof}

\subsection{Sharp flags and Coxeter compatibility}

We next explain why the range containments force the measurements to
be classical at each individual source flag. We do not yet claim
commutation at arbitrary different source labels.

\begin{lemma}[Propagation of incidence words]
\label{lem:transfer-incidence-words}
Consider a finite word of operations, each replacing either one of a
pair of vertices by an incident vertex of a prescribed type, and
ending with dual types. Call the word successful if some sequence of
choices ends at opposite vertices. Success is increasing in the
initial Bruhat position. If a word fails at a source pair, the ordered
candidate product at that pair is supported on target pairs for
which the permuted word fails. The sets of words successful at the
least position are unchanged by the type permutation $p$.
\end{lemma}

\begin{proof}
The base failure rule is nonopposition of dual types, supplied by
Proposition~\ref{prop:transfer-types}; at equal self-opposite vertices
it follows by its continuity argument. Suppose all incident
replacements of one or both source vertices satisfy a previously
proved failure rule. If an incident pair at the target violates that
rule, choose small neighborhoods of the violating vertices.
Proposition~\ref{prop:transfer-incidence} puts the corresponding output
ranges inside sums of the incident source-label ranges. The prior
failure rule makes those product ranges orthogonal. Local homogeneous
sections ensure that nearby base faces still admit incident values
in the chosen neighborhoods. Thus the original ordered product is
supported where every such incident target pair obeys the rule.
Induction on the word proves failure propagation. This argument uses
ordered products and range sums, not a simultaneous Borel calculus
for noncommuting measurements.

For the combinatorics, extend a vertex to a chamber in an apartment.
Changing that chamber without changing the vertex permits its maximal
parabolic Weyl subgroup. An incidence move permits the corresponding
successive maximal-parabolic blocks. The largest attainable Bruhat
positions are obtained by successive ascent to block maxima, that
is, by Demazure multiplication. They are increasing in the initial
position by the Coxeter subword rule, and the apartment realizes the
maxima. Ending opposite at dual vertex types is equivalent to
reaching the longest chamber position after the endpoint extensions.
These are the apartment and Bruhat descriptions of parabolic
incidence~\cite{Borel1991,Tits1964}, and prove monotonicity of success.

If a word is successful at the least target position, it is successful
at every target position. Were its inverse-permuted word to fail at
a least source pair, failure propagation would force the unital
ordered source product to have empty support. Its total product is
the identity, so that is impossible. This gives inclusion of the
least-position success languages in one direction. Iterating $p$,
which has finite order on the type list, gives the cyclic chain of
inclusions back to the original language. They are therefore all
equalities.
\end{proof}

\begin{lemma}[An inequality-detecting word]
\label{lem:transfer-inequality-word}
Every higher-rank component has a vertex type $j$ and an incidence
word that fails exactly when two initial type-$j$ vertices are equal.
\end{lemma}

\begin{proof}
For type $A$, take an end node. Move the second vertex to an incident
vertex of the opposite end type. The word succeeds precisely when a
hyperplane containing the second point can avoid the first point,
which is equivalent to inequality; the apartment point--hyperplane
test gives the same assertion in the building language.

In every other crystallographic diagram choose a self-opposite leaf
$j$. Let $k$ be its neighbor, write $W_{\widehat j}$ for the subgroup
omitting $s_j$, and consider round trips from type $j$ to type $k$
and back. Omitting $k$ isolates $j$, so each round trip permits
exactly the Weyl block
\[
 W_{\widehat j}\{1,s_j\}W_{\widehat j}.
\]
Let $r$ be the minimum number of occurrences of $s_j$ in a reduced
word for the longest element of this component. Starting at equality,
$r-1$ round trips cannot end opposite: an expression so obtained has
at most $r-1$ occurrences of $s_j$, and a reduced subword reaching the
longest element would contradict the definition of $r$.

Starting instead at the least nontrivial double coset
$W_{\widehat j}s_jW_{\widehat j}$, the first such block is already
available. Choose a reduced longest word with exactly $r$ occurrences
of $s_j$ and separate it into the intervening
$W_{\widehat j}$ blocks. The remaining $r-1$ round trips realize it.
Every unequal position is above that least nontrivial double coset,
so success at all unequal pairs follows from
Lemma~\ref{lem:transfer-incidence-words}. The self-opposite leaf exists
in the non-$A$ crystallographic list, including the nonreduced cases
through their underlying Weyl diagram.
\end{proof}

\begin{proposition}[Sharp flag system]
\label{prop:transfer-sharp-flags}
All candidates are sharp and independent of every candidate choice.
For each source vertex $z$, face $F$, and chamber $C$, write the
resulting PVMs as $P_z,P_F,P_C$. Incident source vertex measurements
commute and are supported on incidence of their target values. Every
lexicographic face path, and every mixed face path, towards $F$
observes exactly the vertices of $p(F)$ with these measurements.
Incident source labels also commute across the
two normal measured copies.
\end{proposition}

\begin{proof}
Apply Lemma~\ref{lem:transfer-inequality-word} in a target higher-rank
component, and pull the least-position language back by
Lemma~\ref{lem:transfer-incidence-words}. The inverse-permuted word
fails at equality of source type $p^{-1}(j)$. Failure propagation
therefore puts every ordered product of two candidates at that same
source label on the target diagonal. The unital-POVM argument at the
end of Proposition~\ref{prop:transfer-isolated-sharpness} makes them
equal projections. Together with that proposition, this supplies at
least one sharp type in every target component.

Extend any vertex path lexicographically to a full chamber path. Its
original marginal is retained. If an already sharp vertex label is
incident to that chamber, Proposition~\ref{prop:transfer-incidence}
puts the corresponding full-path marginal ranges inside that sharp
PVM's ranges. Unitality and the same assertion for complements make
the marginal equal to the sharp PVM. A sharp marginal intertwines
the minimal dilation: equality in the CP Schwarz inequality on its
projections gives the intertwining directly. Thus the tested candidate
commutes with the sharp incident measurement, and their ordered
product is supported on target incidence.

Let a remaining target type lie in the component of a sharp type.
For two candidates at the same source label, suppose their target
values $u,v$ have different incidence families of the sharp type.
Choose a sharp-type vertex $w$ incident to $u$ but not to $v$, with
small buffered neighborhoods. The first range containment puts the
ranges near $u$ in the sum of the incident source sharp-label ranges
near $w$. Every one of those ranges is orthogonal to the ranges near
$v$, by the full-path incidence rule just proved. This works with
two independent full-path choices for the tested candidate. Hence
its self-product, and products between choices, are supported where
the two sharp-type incidence families agree.

Those families determine the vertex. Their stabilizer contains its
maximal parabolic and is proper, because the connected apartment
contains nonincident pairs of the relevant types. Bruhat maximality
of the maximal parabolic makes that stabilizer precisely the original
vertex parabolic. Thus equal families imply $u=v$. Diagonal support
again proves sharpness and uniqueness. Repeating in each component
makes all candidates sharp.

We now treat an arbitrary mixed face path of source type $I$.
Fix a positive target stratum with normalized image measurement
$\Psi_F$. For $i\in I$, the type-$i$ source vertex $z_i$ incident
to $F$ is unique. Proposition~\ref{prop:transfer-incidence}\textup{(i)}
puts the support of $\Psi_F(\mathcal I_i(O))$ inside $P_{z_i}(O)$.
A target face omitting $p(i)$ admits two distinct incident vertices
of that type. Choose disjoint basic open patches $O,O'$ around them.
The effect $\Psi_F(\mathcal I_i(O)\cap\mathcal I_i(O'))$ is
supported in both orthogonal projections $P_{z_i}(O),P_{z_i}(O')$,
and hence is zero. Countably many such pairs exhaust the omitted-type
event. Thus every $p(i)$, $i\in I$, occurs. The general upper bound
of Proposition~\ref{prop:links-gap-count} forces the stratum type
to equal $p(I)$. On this stratum $\mathcal I_i(O)$ is precisely
the event that its type-$p(i)$ vertex lies in $O$. The same support
containment and unitality identify its marginal with $P_{z_i}$.
These sharp marginals put the joint face measurement in their
common multiplicative domain, giving $P_F$. This argument uses
mixing, but no order among the diverging coordinates.

For lexicographic paths, the full-path marginals are now uniquely identified and commute by
the sharp-marginal argument. Their joint laws give $P_F,P_C$ and the
stated incidence support. We also verify that a lexicographic path towards $F$ has
no additional observed vertex type. Extend it to a full path and use
persistence and the gap count. An additional sharp marginal would be
invariant under the source face parabolic of $F$ as well as under its
own maximal parabolic. These generate the entire relevant simple
factor, because the extra source type is not retained by $F$.
For a positive subpiece, its scalar mass times the sharp marginal is
its entire angular effect: a POVM dominated by a PVM and having scalar
mass $cI$ equals $c$ times that PVM, by restriction to a projection
and its complement. Normalizing the positive subpiece would therefore
produce a source-invariant boundary measurement. Its finite invariant
trace is forbidden by flag contraction. Thus precisely $p(F)$ is
observed. In the three-slot system, the two normal copies are
marginals of one common full-path dilation. Their sharp incident
marginals therefore commute as well.
\end{proof}

\begin{proposition}[Coxeter compatibility and ordered galleries]
\label{prop:transfer-galleries}
The permutation $p$ preserves the Coxeter matrix. Along a source
chamber gallery, the ordered product of the chamber PVMs is supported
on target galleries of the permuted type, allowing stutters. More
precisely, smooth functions vanishing on that target gallery manifold
annihilate the ordered product, with arbitrary smooth multipliers.
At finite places the corresponding assertion holds for locally
constant tests by compact clopen partitions.
\end{proposition}

\begin{proof}
Adjacent source chambers share their sharp common panel. Insert its
projection between the corresponding consecutive spectral slots.
Differences of the two common-panel coordinates therefore annihilate
the ordered product, including after multiplication by arbitrary
smooth functions expressed by summably separated expansions. The
common-panel equations define the gallery space regularly: it is a
successive fibre product of the submersions from chambers to panels.
They consequently generate its local smooth vanishing ideal, including
at stutters. Smooth partitions of unity prove the asserted
annihilation. In totally disconnected coordinates, finite clopen
partitions insert the common panel projection exactly and give the
same conclusion for locally constant tests.

Fullness gives a nonzero opposite-target restriction of the product
at some generic, hence opposite, source chamber pair. Source covariance
transports this assertion over the opposite orbit. Insert all
intermediate PVM identities along any reduced longest source gallery,
and use compact partitions of the intermediate target chamber spaces.
Nonzero opposite endpoint support requires a gallery of the permuted
type connecting opposite target chambers. Both ends use the full
spherical Weyl system of the fixed ambient group $G$. Their longest
elements therefore have the same Coxeter length. The permuted word must
therefore itself be a reduced
longest word; no stutter or nonminimal step is possible.

An alternating rank-two word of Coxeter length $m_{ij}$ is reduced
and extends to a reduced longest word. Its permuted segment is also
reduced, so $m_{ij}\leq m_{p(i)p(j)}$. Cycling the finite permutation
$p$ gives equality. This proves Coxeter compatibility.
\end{proof}

\subsection{Exact residue links}

Take a mixed lex face path missing one type in a higher-rank
component, or one type in each of several higher-rank components.
Fix one such component. Its source rank-one Levi boundary parametrizes
matched opposite-panel completions at the two source ends. Denote
this common Levi flag by $z$. Refining each of the two opposite
source faces by that same flag of their common Levi gives its
start and ordinary-end missing vertices. This is the transverse
parabolic identification of the residues, with the ordinary-end
types relabeled by opposition. Write $P_z$ and $Q_z$ for the sharp
missing-vertex PVMs at these start and ordinary-end vertices,
respectively. Their common source parameter is transported by
$z\mapsto sz$ under the simultaneous Levi action. Frame the two sharp endpoint
faces separately by compact transports, and identify the reference
opposite panels by transverse projection. The residual output
coordinate $g$ is then a rank-one target projectivity from the end
pencil to the start pencil.

\begin{proposition}[Residue graph equation]
\label{prop:transfer-residue-graph}
On compact output windows the mixed link satisfies
\begin{equation}
 E(dg)J\bigl(\overline{Q_z}(dq)\otimes P_z(du)\bigr)
 \quad\text{is supported on}\quad u=gq.
 \label{eq:input-6}
\end{equation}
The equation means annihilation by smooth joint tests vanishing on
the graph, with locally constant tests in finite-place angular
variables; conjugate spectral measures conjugate their projections.
It holds with bounded Borel output parameters after the angular
identities have been established. The joint law is invariant under
simultaneous source Levi transport, and exact input translations may
be made before any further limit.
\end{proposition}

\begin{proof}
We first work with an archimedean target component. There are two
steps, both using ordered spectral slots.

\emph{Projection onto the reverse pencil.}
On the ordinary-end Hilbert space, use the slots in the order:
reverse panel, end missing vertex at $z$, and the observation at the
start label, possibly moved by the reverse radical. Denote the value
in the last slot by $v$. A source gallery shortened from the longest
by the missing reflection connects the matched chambers. Its target
gallery has the same shortened type by
Proposition~\ref{prop:transfer-galleries}. The shortened Weyl position
is $s_iw_0=w_0s_{i^*}$; it has all the other respective left and right
descents, so saturation by the parabolics omitting those missing
reflections cannot reach $w_0$. Thus the missing vertices are
nonopposite.

A vertex transverse to a panel of the dual missing type has exactly
one incident completion of that panel whose missing vertex is
nonopposite to it. This is the ordinary rank-one parabolic projection.
The ordered gallery equation therefore says that the end missing
vertex is exactly $\operatorname{pr}(v)$ on a transverse window.
Integrating the intermediate gallery slots by their identity
projections preserves the equation. Other components can be completed
arbitrarily.

\emph{The line equation.}
Use \eqref{eq:input-5} at the single start vertex $z$ and move only its
end occurrence by the independent reverse radical. Sharpness permits
taking the start column to be the projector
$1-n(P_z)n(P_z)^*$ applied to an arbitrary column, and the end column
to be the moved line projector on a transverse patch. Thus, after
bounded output localization, the link is annihilated by
\[
 (1-n(u)n(u)^*)S(h)n(v).
\]
Its two angular inputs commute because they act on the two separate
input tensor slots. The identity consequently holds with arbitrary
smooth joint angular multipliers, by summable separation. On the
stratum, the matrix action is the residual projectivity after
transverse projection, so the equations are precisely the line graph
$u=g\operatorname{pr}(v)$. They generate its smooth vanishing ideal
on compact transverse windows: the highest-line embedding is
immersive, so the derivative in the image-vertex argument has full
injective rank. Smooth division is uniform on such windows. Output
coordinates may be Borel parameters with locally uniform angular
derivative bounds, since separation is performed only in the angular
variables.

We now combine the two steps in their necessary order. Preapply at
the end a reverse-panel cut followed by a moved-vertex cut, on a
buffered transverse rectangle, and leave the start vector arbitrary.
Insert the $Q_z$ spectral slot next to the reverse-panel slot; they
commute. The shortened-gallery equation replaces its value by
$\operatorname{pr}(v)$, including after smooth multiplication in the
independent start and output variables. Integrate out that slot.
Pull the sharp reverse-panel projection through $J$, using
Proposition~\ref{prop:links-link}. The line equation then supplies
$u=g\operatorname{pr}(v)=gq$. This proves
\eqref{eq:input-6} on those ordered ranges. In particular, we have not
applied a discontinuous joint symbol to an arbitrary pair of
noncommuting end measurements.

These ordered ranges are total. Their closed join is target
invariant. It is fixed by a split source subgroup escaping in every
ambient component, fixing the two panels and their respective
completions, and normalizing the reverse-radical moves. Its projection
is therefore scalar by the invariant-mass and mixing rule of
Proposition~\ref{prop:links-link}. Some such ranges are nonzero when
also varying $z$: moved start labels then cover generic vertices, and
patch fullness applies to a nonzero reverse-panel projection on a
uniform transverse target patch. Source Levi transitivity on the
completion parameter transports nonzeroness to every $z$. Hence the
join is the whole end space, and the equation holds everywhere.

\emph{Finite target places.}
At a finite target place, replace the Hilbert line argument by the
individual-pair assertion of
Proposition~\ref{prop:transfer-incidence}. Choose scalar clopen start
and end patches. On a buffered output window carrying the end patch
strictly inside the start image patch, a start vector in the kernel
of the corresponding $P_z$ projection and an end vector in the cut
range of the same start-label measurement give zero. Since $P_z$ is
sharp, finite clopen partitions imply the exact transverse image
substitution. Move the end occurrence by the reverse radical as
before. The action on transverse values is the residual projectivity
after projection to the reverse pencil.

The shortened-gallery substitution holds in this component by compact
clopen partitions and insertion of the common sharp panel projection.
Thus the reverse-panel cut followed by the moved-vertex cut again
imposes $u=gq$, including arbitrary locally constant angular
multipliers. The same total-range argument applies. This proves the
finite-place equation without tensoring the complex Hilbert input
with a representation over a nonarchimedean field.

For mixed target coordinates, use finite partitions in the clopen
slots and summable smooth separated expansions in the real slots.
Individual stratum restrictions are made by Borel output projections
after buffered tests in the ambient compactification; no uniform
transverse estimate for an arbitrary discontinuous mask in an
unbounded-overlap problem is required.

Finally, a source Levi element centralizes the chosen face path.
Exact cocycle composition, made before limiting, shows that its
simultaneous action at the two inputs changes the array only by base
reindexing and a common scalar phase. The tested forms are therefore
invariant. If a genuine action is needed, use the conjugate-multiplier
regular amplification from Section~\ref{sec:links}. It tensors the
forms by identity and cancels the phase defect at each input, without
altering any finite-trace mixing assertion. This proves the stated
transport properties.
\end{proof}

\begin{corollary}[Fibrewise links and polynomial fullness]
\label{cor:transfer-fibre-links}
Disintegrate over the two sharp endpoint faces and use genuine source
amplifications if necessary. For product-almost every pair of endpoint
fibres, the mixed isometry \eqref{eq:input-1}, the graph equation
\eqref{eq:input-6}, simultaneous source Levi invariance, and exact
input translations hold. On almost every individual fibre, for every
even highest power in an archimedean target module, the line ranges
over $z$ span at least the Hilbert fibre tensored with the span of all
ordinary incident highest lines.

Consequently, a homogeneous polynomial relation with linear
Hilbert-vector coefficients that holds spectrally at every residue
label holds at every ordinary residue value. The assertion includes
algebraic subspace membership expressed by Pluecker equations and
chartwise relations after the denominator procedure in the proof.
\end{corollary}

\begin{proof}
The sharp face marginals intertwine the link dilation, so its spectral
span disintegrates over the product of the two endpoint spectral
measures. The strongly continuous source Levi representations commute
with those face PVMs and therefore disintegrate as well. One may obtain
continuous fibre representations by disintegrating a countable dense
convolution algebra of the group and then using its nondegenerate
integrated representations, as in the standard direct-integral
construction~\cite[Theorem~1.G.6]{BekkaHarpe2019}.

Before removing a null set, choose countably many dense residue
labels, compact windows, local graph generators, angular Fourier
tests in real charts, fixed-level clopen tests in finite charts, and
total input vectors. All the corresponding identities then hold on
one conull set. Strong source transport and the uniform angular
bounds extend them to every residue label and smooth or locally
constant angular test there. The stabilizer invariance at a standard
label also disintegrates, first on a countable dense subgroup and
then by strong continuity, so the transported family is well defined
on the residue quotient. Once a vector identity holds on an individual
fibre, any bounded Borel output multiplier is applied directly through
that fibre's spectral calculus. Uniformly summable angular expansions
then give tests with Borel output parameters. No countable norm-dense
class of all bounded Borel functions is used.

Disintegrate the amplified incidence containment of
Proposition~\ref{prop:transfer-incidence} on the same determining
family. On almost every fibre its closed span contains the Hilbert
fibre tensored with the span of all ordinary incident highest lines,
and this holds separately for every even highest power. A polynomial
with linear Hilbert-vector coefficients that vanishes on all spectral
line ranges is orthogonal to their span. It is therefore orthogonal
to each ordinary incident highest line, proving polynomial fullness.
Norm powers homogenize different even degrees.

For a relation known only on a chart $\{\Delta\ne0\}$, first clear a
common homogeneous denominator and call the resulting numerator
$F_{\mathrm{num}}$. Apply global polynomial fullness to
$\Delta F_{\mathrm{num}}$: it vanishes spectrally on the chart by the
relation, and on $\{\Delta=0\}$ by the extra factor. Divide by
$\Delta$ on the original chart to recover the relation there.
Pluecker equations and common local denominators put the stated
algebraic membership tests in this form.
\end{proof}
